% =============================================================================
% Section 6: DAI: Timing Side-Channel Defense
% =============================================================================

\section{DAI: Timing Side-Channel Defense}
\label{sec:dai}

Theorem~\ref{thm:isolation} guarantees that no private data leaks through the attention
\emph{data channel}. However, the shared KV cache introduces a potential \emph{metadata
channel}: cache hit/miss timing may reveal whether a page is resident, leaking occupancy
patterns~\cite{osvik2006cache, yarom2014flush, percival2005cache}. We introduce the
\textbf{Deterministic Anomaly Interceptor (DAI)}, a lightweight per-tenant state machine
that detects and mitigates timing side-channel attacks in real time.

% -----------------------------------------------------------------------------
\subsection{Collision Attack Model}
\label{sec:dai:attack}

We adopt the standard cache timing oracle model~\cite{osvik2006cache, percival2005cache}.
An adversary controlling tenant $t_a$ issues carefully crafted queries to probe whether
a target tenant $t_v$'s pages are resident in the KV cache. The observed latency for a
single probe is:
\begin{equation}
    T_{\mathrm{obs}} = \begin{cases}
        T_{\mathrm{hit}} + \eta & \text{if the target page is cached}, \\
        T_{\mathrm{miss}} + \eta & \text{if the target page is evicted},
    \end{cases}
    \label{eq:timing-oracle}
\end{equation}
where $T_{\mathrm{hit}} < T_{\mathrm{miss}}$ are the deterministic cache hit and miss
latencies, $\eta \sim \mathcal{N}(0, \sigma_{\mathrm{sys}}^2)$ is system noise, and
$\Delta \triangleq T_{\mathrm{miss}} - T_{\mathrm{hit}}$ is the timing gap.

Let $S_{\mathrm{cache}} \in \{0, 1\}$ denote the binary cache state (0 = miss, 1 = hit).
The adversary's observation forms a \emph{binary channel}~\cite{shannon1948mathematical}
with capacity:
\begin{equation}
    C = 1 - H_b(p_e),
    \label{eq:channel-capacity}
\end{equation}
where $H_b(p_e) = -p_e \log_2 p_e - (1 - p_e)\log_2(1 - p_e)$ is the binary entropy
function and $p_e = \Phi(-\Delta / (2\sigma_{\mathrm{sys}}))$ is the crossover probability
under optimal thresholding ($\Phi$ is the standard normal CDF).

\textbf{Without mitigation}, the adversary can average over $n$ independent probes,
reducing the effective noise to $\sigma_{\mathrm{eff}} = \sigma_{\mathrm{sys}} / \sqrt{n}$.
As $n \to \infty$, $p_e \to 0$, yielding $C \to 1$ bit per probe---a perfect timing oracle.

% -----------------------------------------------------------------------------
\subsection{DFA-Based Detection}
\label{sec:dai:dfa}

DAI models each tenant's behavior as a \textbf{deterministic finite automaton (DFA)} with
four states, enabling stateful detection that distinguishes transient anomalies from
sustained attack patterns.

\begin{definition}[DAI State Machine]
The DFA is the tuple $(\mathcal{S}, \Sigma, \delta, s_0, \mathcal{F})$ where:
\begin{itemize}[leftmargin=1.5em,itemsep=2pt]
    \item $\mathcal{S} = \{S_0, S_1, S_2, S_3\}$: states \textsc{Normal},
          \textsc{Probe\_1}, \textsc{Probe\_2}, \textsc{Attack\_Confirmed};
    \item $\Sigma$: per-query statistics (rolling variance $\hat{\sigma}$, query count,
          distributional divergence);
    \item $\delta: \mathcal{S} \times \Sigma \to \mathcal{S}$: transition function
          (Table~\ref{tab:dfa-transitions});
    \item $s_0 = S_0$: initial state;
    \item $\mathcal{F} = \{S_3\}$: the accepting (defense-active) state.
\end{itemize}
\end{definition}

\noindent The transition guards are parameterized by three thresholds:

\begin{enumerate}[leftmargin=1.5em,itemsep=2pt]
    \item \textbf{$S_0 \to S_1$ (anomaly detection):} $\hat{\sigma} > \tau_\sigma$, where
          $\hat{\sigma}$ is the rolling standard deviation of inter-query arrival times
          over a sliding window of $W = 50$ queries, and $\tau_\sigma = 15$~ms.
    \item \textbf{$S_1 \to S_2$ (sustained anomaly):} the condition $\hat{\sigma} > \tau_\sigma$
          persists for $K = 20$ consecutive queries, ruling out transient load spikes.
    \item \textbf{$S_2 \to S_3$ (distributional confirmation):} the KL divergence between
          the observed timing distribution $\hat{P}$ and the benign baseline $P_{\mathrm{benign}}$
          exceeds the threshold:
          $D_{\mathrm{KL}}(\hat{P} \| P_{\mathrm{benign}}) > \kappa$, with $\kappa = 2.0$~nats.
    \item \textbf{$S_3 \to S_0$ (cooldown):} the tenant must satisfy \emph{all three}
          conditions simultaneously for at least $C_{\mathrm{cool}} = 100$ consecutive
          queries: $\hat{\sigma} \leq \tau_\sigma$, $D_{\mathrm{KL}}(\hat{P} \| P_{\mathrm{benign}}) \leq \kappa$,
          and the query rate returns to within $2\times$ of the tenant's historical baseline.
\end{enumerate}

\begin{table}[t]
\centering
\caption{DAI state transition table. Each row specifies the current state, guard condition,
next state, and triggered action.}
\label{tab:dfa-transitions}
\small
\begin{tabular}{@{}ccll@{}}
\toprule
\textbf{From} & \textbf{To} & \textbf{Guard} & \textbf{Action} \\
\midrule
$S_0$ & $S_1$ & $\hat{\sigma} > \tau_\sigma$ & Begin monitoring window \\
$S_1$ & $S_0$ & $\hat{\sigma} \leq \tau_\sigma$ before $K$ queries & Reset counter \\
$S_1$ & $S_2$ & $\hat{\sigma} > \tau_\sigma$ sustained for $K = 20$ queries & Compute $D_{\mathrm{KL}}$ \\
$S_2$ & $S_0$ & $D_{\mathrm{KL}}(\hat{P} \| P_{\mathrm{benign}}) \leq \kappa$ & False alarm; reset \\
$S_2$ & $S_3$ & $D_{\mathrm{KL}}(\hat{P} \| P_{\mathrm{benign}}) > \kappa = 2.0$ nats & Activate countermeasures \\
$S_3$ & $S_3$ & Any guard violated during cooldown & Maintain defense \\
$S_3$ & $S_0$ & $C_{\mathrm{cool}} = 100$ clean queries & Deactivate; resume normal \\
\bottomrule
\end{tabular}
\end{table}

\begin{algorithm}[t]
\caption{DAI Per-Tenant State Machine}
\label{alg:dai}
\begin{algorithmic}[1]
\Require Tenant state $s \in \{S_0, S_1, S_2, S_3\}$, timing buffer $\mathcal{B}$, counter $c$
\Ensure Updated state $s'$, response latency modifier $\delta_t$
\State $\mathcal{B}.\text{append}(T_{\mathrm{current}})$ \Comment{Record query timestamp}
\State $\hat{\sigma} \gets \text{RollingStd}(\mathcal{B}, W{=}50)$
\If{$s = S_0$}
    \If{$\hat{\sigma} > \tau_\sigma$}
        $s' \gets S_1$; $c \gets 1$
    \Else~$s' \gets S_0$
    \EndIf
\ElsIf{$s = S_1$}
    \If{$\hat{\sigma} \leq \tau_\sigma$}
        $s' \gets S_0$; $c \gets 0$ \Comment{Transient; reset}
    \ElsIf{$c \geq K$}
        \State $\hat{P} \gets \text{HistogramEstimate}(\mathcal{B})$
        \If{$D_{\mathrm{KL}}(\hat{P} \| P_{\mathrm{benign}}) > \kappa$}
            $s' \gets S_3$ \Comment{Skip $S_2$: confirmed}
        \Else~$s' \gets S_0$; $c \gets 0$ \Comment{Benign anomaly}
        \EndIf
    \Else~$c \gets c + 1$; $s' \gets S_1$
    \EndIf
\ElsIf{$s = S_3$}
    \State $\delta_t \gets |\mathcal{N}(0, \sigma_{\mathrm{noise}}^2)|$ \Comment{Inject noise}
    \State Apply rate limiting: delay by factor $\lambda = 2.0$
    \If{$\hat{\sigma} \leq \tau_\sigma ~\wedge~ D_{\mathrm{KL}} \leq \kappa ~\wedge~ c_{\mathrm{cool}} \geq C_{\mathrm{cool}}$}
        $s' \gets S_0$ \Comment{Cooldown complete}
    \Else~$s' \gets S_3$; update $c_{\mathrm{cool}}$
    \EndIf
\EndIf
\State \Return $s', \delta_t$
\end{algorithmic}
\end{algorithm}

% -----------------------------------------------------------------------------
\subsection{Countermeasures}
\label{sec:dai:countermeasures}

Upon entering state $S_3$, DAI activates two complementary countermeasures that jointly
destroy the timing channel.

\paragraph{Noise injection.}
Each response to the flagged tenant is delayed by an additive noise term
$\delta_t = |\mathcal{N}(0, \sigma_{\mathrm{noise}}^2)|$ with
$\sigma_{\mathrm{noise}} = 5$~ms (half-normal distribution, ensuring non-negative delay).
This noise is \emph{independent} of the cache state $S_{\mathrm{cache}}$ and is injected
at the serving layer, after attention computation.

\paragraph{Rate limiting.}
The flagged tenant's query throughput is throttled by a factor of $\lambda = 2.0$,
doubling the minimum inter-query interval. This limits the number of samples available
for averaging attacks.

\paragraph{SNR analysis.}
The pre-defense signal-to-noise ratio for a single probe is:
\begin{equation}
    \mathrm{SNR}_{\mathrm{pre}} = \frac{\Delta^2}{\sigma_{\mathrm{sys}}^2}.
\end{equation}
After noise injection, the effective noise variance becomes
$\sigma_{\mathrm{eff}}^2 = \sigma_{\mathrm{sys}}^2 + \sigma_{\mathrm{noise}}^2$, yielding:
\begin{equation}
    \mathrm{SNR}_{\mathrm{post}} = \frac{\Delta^2}{\sigma_{\mathrm{sys}}^2 + \sigma_{\mathrm{noise}}^2}
    \ll \mathrm{SNR}_{\mathrm{pre}}.
    \label{eq:snr-post}
\end{equation}
For typical values ($\Delta = 2$~ms, $\sigma_{\mathrm{sys}} = 1$~ms,
$\sigma_{\mathrm{noise}} = 5$~ms), we have
$\mathrm{SNR}_{\mathrm{post}} = 4 / 26 \approx 0.15$, a $17\times$ reduction.

\paragraph{Effective channel capacity.}
Combining noise injection with rate limiting, the effective channel capacity becomes:
\begin{equation}
    C_{\mathrm{eff}} = \frac{1}{\lambda}\big(1 - H_b(p_e')\big),
    \label{eq:effective-capacity}
\end{equation}
where $p_e' = \Phi\!\big(-\Delta / (2\sqrt{\sigma_{\mathrm{sys}}^2 + \sigma_{\mathrm{noise}}^2})\big)$.
Substituting:
\begin{equation}
    p_e' = \Phi\!\bigg(\frac{-2}{2\sqrt{26}}\bigg) = \Phi(-0.196) \approx 0.422.
\end{equation}
The binary entropy evaluates to $H_b(0.422) \approx 0.982$ bits, giving:
\begin{equation}
    C_{\mathrm{eff}} = \frac{1}{2}(1 - 0.982) = 0.009~\text{bits/probe}.
\end{equation}
This represents a $>99\%$ reduction in channel capacity compared to the undefended
$C \to 1$ bit/probe. Furthermore, rate limiting halves the probe rate, so the
\emph{throughput} of leaked information drops by an additional factor of $2$, yielding
an effective information rate of $\sim$0.0045 bits per unit time---practically zero.

% -----------------------------------------------------------------------------
\subsection{Timing Information Bound}
\label{sec:dai:bound}

\begin{theorem}[Timing Information Bound]
\label{thm:timing-bound}
Under DAI, the mutual information between the cache state and observed timing is bounded:
\begin{equation}
    I(S_{\mathrm{cache}};\; T_{\mathrm{obs}}) \leq I_{\mathrm{natural}} + \varepsilon_{\mathrm{FNR}},
\end{equation}
where $I_{\mathrm{natural}}$ is the mutual information under benign operation
(due to inherent system timing variation) and $\varepsilon_{\mathrm{FNR}}$ is a residual
term bounded by the false negative rate of the DFA detector.
\end{theorem}

\begin{proof}
We analyze three temporal phases of any attack.

\medskip
\noindent\textbf{Phase 1: Pre-detection (states $S_0 \to S_1 \to S_2$).}
Before DAI transitions to $S_3$, the adversary operates with the undefended channel.
The maximum number of queries in this phase is bounded by the detection latency:
$n_{\mathrm{pre}} \leq K + W = 20 + 50 = 70$ queries (the monitoring window $W$ plus
the sustained anomaly threshold $K$). During this phase, the adversary accumulates at most:
\begin{equation}
    I_{\mathrm{pre}} \leq n_{\mathrm{pre}} \cdot C \leq 70 \cdot 1 = 70~\text{bits}
\end{equation}
of timing information. This is a \emph{finite, bounded} quantity, independent of the
total attack duration.

\medskip
\noindent\textbf{Phase 2: Active defense (state $S_3$).}
Once countermeasures activate, the observed timing becomes:
\begin{equation}
    T_{\mathrm{obs}} = T_{\mathrm{true}} + \delta_t, \quad \delta_t \sim |\mathcal{N}(0, \sigma_{\mathrm{noise}}^2)|,
    \quad \delta_t \perp\!\!\!\perp S_{\mathrm{cache}}.
\end{equation}
By the \textbf{data processing inequality}~\cite{cover2006elements}:
\begin{equation}
    I(S_{\mathrm{cache}}; T_{\mathrm{obs}}) \leq I(S_{\mathrm{cache}}; T_{\mathrm{true}}).
\end{equation}
But $T_{\mathrm{obs}}$ is a noisy function of $T_{\mathrm{true}}$, and the additive
independent noise strictly reduces mutual information. Specifically, by
Equation~\eqref{eq:effective-capacity}, the per-probe capacity drops to
$C_{\mathrm{eff}} \approx 0.009$ bits. Combined with rate limiting ($\lambda = 2.0$),
the information rate is:
\begin{equation}
    R_{\mathrm{attack}} = \frac{C_{\mathrm{eff}}}{\lambda} \approx 0.0045~\text{bits/query} \to 0
    \quad \text{as } \sigma_{\mathrm{noise}} / \Delta \to \infty.
\end{equation}

\medskip
\noindent\textbf{Phase 3: Steady state.}
Under a persistent attack ($\hat{\sigma}$ remains elevated), the DFA remains in $S_3$
indefinitely. The cooldown guard ($C_{\mathrm{cool}} = 100$ consecutive clean queries
with $\hat{\sigma} \leq \tau_\sigma$ \emph{and} $D_{\mathrm{KL}} \leq \kappa$) prevents
premature relaxation. If the attacker ceases probing to satisfy the cooldown condition,
they are by definition no longer attacking, and the channel is idle.

If the attacker resumes after cooldown, the DFA re-enters $S_1$ within one anomalous
query and reaches $S_3$ within $K = 20$ queries, bounding the new pre-detection leakage
to another $\leq 70$ bits---a pattern that yields at most $70$ bits per attack epoch,
with epoch length growing due to the mandatory $100$-query cooldown.

\medskip
\noindent\textbf{Combining the phases.}
The total steady-state leakage rate is:
\begin{equation}
    I(S_{\mathrm{cache}}; T_{\mathrm{obs}}) \leq I_{\mathrm{natural}} + \varepsilon_{\mathrm{FNR}},
\end{equation}
where $I_{\mathrm{natural}}$ captures the unavoidable timing variation of benign cache
behavior (present in \emph{any} caching system) and
$\varepsilon_{\mathrm{FNR}} = P(\text{FN}) \cdot C$ accounts for the residual leakage
during any undetected attack queries. In our empirical evaluation (\S\ref{sec:eval}),
the DFA achieves a false negative rate below $10^{-3}$, yielding
$\varepsilon_{\mathrm{FNR}} < 10^{-3}$ bits/query.
\end{proof}

\begin{remark}[Defense in Depth]
\label{rem:defense-in-depth}
TSAM and DAI provide complementary, layered protection~\cite{denning1976lattice, goguen1982security}:
\begin{enumerate}[leftmargin=1.5em,itemsep=2pt]
    \item \textbf{Data channel} (TSAM): Theorem~\ref{thm:isolation} guarantees
          $I(\mathrm{Attn}_{t_i}; \mathrm{KV}_{\mathrm{private}}(t_j)) = 0$ exactly---the
          attention mechanism reveals \emph{zero bits} of private data.
    \item \textbf{Metadata channel} (DAI): Theorem~\ref{thm:timing-bound} bounds the
          timing side-channel to $I_{\mathrm{natural}} + \varepsilon_{\mathrm{FNR}}
          \approx 0$---the cache timing reveals \emph{negligible bits} of occupancy
          metadata.
\end{enumerate}
Together, these two mechanisms close both the primary data exfiltration vector and the
secondary timing side-channel, achieving comprehensive tenant isolation without modifying
model weights, degrading output quality, or imposing significant computational overhead.
\end{remark}
